\chapter{Background and Definition}\label{ch:background}\
Before exploring the literature on lifelong learning techniques, it is important to set appropriate definitions of several widely used terms, provide brief explanations of key concepts, and to give sufficient context for the scope of the research.

\section{Lifelong, Continual, or Incremental Learning?}\label{sec:lifelong-def}\
The broad concept of lifelong learning has been given many names throughout the literature.
Common terms are (among others) continual, incremental, and online learning.\
There are some nuances between these terminologies when it comes to the exact details of how regularly models are updated (retrained), the batch sizes of the update training sets, and whether the data shifts over time.
`Lifelong' simply implies maintaining good functionality throughout the lifetime, `continual' includes shifting data distributions, `incremental' suggests at least regular updates with potentially single training samples, and `online' refers to updates during test time.
However, in this study, the terms above will be used interchangeably.\ Moreover, `lifelong' will be used as a broad overarching term encompassing all the aspects above.
Notably, using these terms, no distinction is made on whether they convey real-time implementations.


\section{Learning vs. Adapting}\label{sec:learning-v-adapting}\
Throughout \autoref{ch:introduction}, it was mentioned several times that the goal is to have a robot continually
adapt to changes.
Apparently, `adapting' and `learning' can be used interchangeably in this context, or could (and will be throughout this study) both be interpreted as `learning to adapt'.
However, some distinction between the two can be made to provide some context.\
% important definitions: continual vs incremental vs lifelong; learning vs adapting;

% historical background

% ML processes: RL, DP, etc (?); model-based, policy-based, ...

% learning paradigms